\chapter{2022年「大学物理2」\hdu 补考试题}
\date{2023年3月14日}{大学物理教学团队}{A0715011}
{未央物理讲师 \href{https://qm.qq.com/q/UPbGudx8cK}{\textbf{Axia}}}
{\qrcode{https://qm.qq.com/q/asJiytD5pC}}{「未央」HDU 物理营}
{\qrcode{http://weixin.qq.com/r/mhKYgFTEMQlOrRDj90eI}}{未央学社公众号}
% \input{cha/guide22DW-B.tex}
\section{单选题（共30分，每小题3分）}
{
\begin{choice}[B]{简谐振动}
    一质点作简谐振动，振动方程为$x=A\cos{\ab(\omega t+\varphi)}$，当时间$t=\frac{T}{2}$（$T$为周期）时，质点的速度为
    \begin{tasks}(4)
        \task $-\omega A\sin{\varphi}$
        \task $\omega A\sin{\varphi}$
        \task $-\omega A\cos{\varphi}$
        \task $\omega A\cos{\varphi}$
    \end{tasks}
\end{choice}
\solute{$v\ab(T/2)=\eval{\frac{\d x}{\d t}}_{t=T/2}=-\omega A\sin{\ab(\omega\frac{T}{2}+\varphi)}\xlongequal{T=\frac{2\pi}{\omega}}\omega A\sin{\varphi}$.\sokka{B}}

\begin{choice}[D]{简谐振动}
    一弹簧振子作简谐振动，总能量为$E_1$，如果简谐振动振幅增加为原来的两倍，重物的质量增为原来的$4$倍，则它的总能量$E_2$变为
    \begin{tasks}(4)
        \task $\frac{E_1}{4}$
        \task $\frac{E_1}{2}$
        \task $2E_1$
        \task $4E_1$
    \end{tasks}
\end{choice}
\solute{弹簧振子总能量$E=\frac{1}{2}kA^2$与振幅有关，与重物质量无关.所以当振幅增加为原来的两倍时，总能量变为原来的$4$倍.\sokka{D}}

\begin{choice}[C]{平面简谐波的物理量}
    一平面简谐波表达式为$y=-0.05\sin{\pi\ab(t-2x)}$ (SI).则该波的频率$\nu$，波速$u$及波线上各点振动的振幅$A$依次为
    \begin{tasks}(4)
        \task $\frac{1}{2},\ -\frac{1}{2},\ -0.05$
        \task $\frac{1}{2},\ 1,\ -0.05$
        \task $\frac{1}{2},\ -\frac{1}{2},\ 0.05$
        \task $2,\ 2,\ ,-0.05$
    \end{tasks}
\end{choice}
\solute{题中的波函数可改写为$y=0.05\sin{\ab[2\pi\ab(\frac{1}{2}t-x)-\frac{\pi}{2}]}=A\cos{\ab[2\pi\ab(\nu t-\frac{x}{\lambda})+\varphi]}$，由此可得$\nu=\frac{1}{2},\ \lambda=1,\ u=\lambda\nu=\frac{1}{2},\ A=0.05$.\sokka{C}}

\begin{choice}[C]{}
    在同一媒质中两列相干的平面简谐波的强度之比是$\frac{I_1}{I_2}=4$，则两列波的振幅之比是
    \begin{tasks}(4)
        \task $\frac{A_1}{A_2}=16$
        \task $\frac{A_1}{A_2}=4$
        \task $\frac{A_1}{A_2}=2$
        \task $\frac{A_1}{A_2}=\frac{1}{4}$
    \end{tasks}
\end{choice}
\solute{平面简谐波的强度之比与振幅之比的平方成正比，故$\frac{A_1}{A_2}=\sqrt{\frac{I_1}{I_x}}=2$.\sokka{C}}

\begin{choice}[B]{}
    在波长为$\lambda$的驻波中，两个相邻波腹之间的距离为
    \begin{tasks}(4)
        \task $\frac{\lambda}{4}$
        \task $\frac{\lambda}{2}$
        \task $\frac{3\lambda}{4}$
        \task $\lambda$
    \end{tasks}
\end{choice}

\begin{choice}[A]{双缝干涉}
    薄钢片上有两条紧靠着的平行细缝，用双缝干涉方法来测量两缝间距.如果用波长$\lambda=546.1\mathrm{nm}$的单色光照射，双缝与屏的距离$D=300\mathrm{mm}$，测得中央明条纹两侧的两个第五级明条纹的间距为$12.2\mathrm{mm}$，则两缝间距离为
    \begin{tasks}(4)
        \task $0.134\mathrm{mm}$
        \task $1.22\mathrm{mm}$
        \task $1.34\mathrm{mm}$
        \task $0.122\mathrm{mm}$
    \end{tasks}
\end{choice}
\solute{由明纹位置$x=k\frac{\lambda D}{d}$得两个第五级明条纹的间距$\Delta x_{\pm 5}=5\frac{\lambda D}{d}-\ab(-5\frac{\lambda D}{d})=10\frac{\lambda D}{d}=12.2\mathrm{mm}$.代入$D,\ \lambda$的数据得$d=0.134\mathrm{mm}$.\sokka{A}}

\begin{choice}[C]{}
    当牛顿环干涉装置的透镜垂直于平板玻璃向上平移（离开玻璃板时），从入射光方向观察到干涉环纹的变化情况是
    \begin{tasks}(2)
        \task 环纹向边缘扩散，环数不变
        \task 环纹向边缘扩散，环数增加
        \task 环纹向中心扩散，环数不变
        \task 环纹向中心扩散，环数减少
    \end{tasks}
\end{choice}
\solute{当透镜离开玻璃板时，两者之间的空气薄膜厚度增加，光程差增加；而牛顿环的条纹级数与光程差成正比（意思就是某处原本是$n$级别条纹，被这么一操作，变成了$n+k$级，原本的$n$级条纹跑到了里面），所以环纹向中心扩散；而环纹间距与此无关，所以环纹级数不变.\sokka{C}}

\begin{choice}[B]{钟慢效应}
    在某地发生两件事，静止位于该地的甲测得时间间隔为$4\mathrm{s}$，若相对于甲作匀速直线运动的乙测得时间间隔为$5\mathrm{s}$，则乙相对于甲的运动速度是（$c$表示真空中光速）
    \begin{tasks}(4)
        \task $\frac{4}{5}c$
        \task $\frac{3}{5}c$
        \task $\frac{2}{5}c$
        \task $\frac{1}{5}c$
    \end{tasks}
\end{choice}
\solute{由钟慢效应得乙所测时间$t=\frac{t_0}{\sqrt{1-v^2/c^2}}=5\mathrm{s}$，所以乙相对于甲的速度为$v=0.6c$.\sokka{B}}

\begin{choice}[A]{}
    粒子动能等于它本身的静止能量，这时粒子的速率为
    \begin{tasks}(4)
        \task $\frac{\sqrt{3}}{2}c$
        \task $\frac{3}{4}c$
        \task $\frac{1}{2}c$
        \task $\frac{4}{5}c$
    \end{tasks}
\end{choice}
\solute{由题意得$E_k=mc^2-m_0c^2=m_0c^2$，根据质速关系$m=\frac{m_0}{\sqrt{1-v^2/c^2}}$得$v=\frac{\sqrt{3}}{2}c$.\sokka{A}}

\begin{choice}[A]{}
    在氢原子的$K$壳层中，电子可能具有的量子数$(n,\ l,\ m_l,\ m_s)$是
    \begin{tasks}(4)
        \task $(1,\ 0,\ 0,\ \frac{1}{2})$
        \task $(1,\ 0,\ -1,\ \frac{1}{2})$
        \task $(1,\ 1,\ 0,\ -\frac{1}{2})$
        \task $(2,\ 1,\ 0,\ -\frac{1}{2})$
    \end{tasks}
\end{choice}
\solute{四个量子数关系：$n$确定时，$l=0,\pm 1,\cdots,\pm n-1;\ m_l=0,\pm 1,\cdots,\pm l;\ m_s=\pm\frac{1}{2}$.\sokka{A}}
}

\section{填空题（共22分）}
\begin{exercise}[2]{双缝干涉}
    若将双缝干涉实验从空气移入水面之下进行，则干涉条纹间的距离将\ans{变小}（填变大、变小或不变）.
\end{exercise}
\solute{移入水中后明纹的光程差$\delta=\frac{k\lambda}{n}=\frac{nxd}{D}$，由此得条纹间距$\Delta x=\frac{k\lambda D}{n^2d}$.由于$n>1$，所以间距变小.}

\begin{exercise}[2]{双缝干涉}
单色平行光垂直入射到双缝上，观察屏上$P$点到两缝的距离分别为$r_1$和$r_2$.设双缝和屏之间充满折射率为$n$的妹纸，则$P$点处两相干光波的光程差为\ans{$n\ab(r_2-r_1)$}.
\end{exercise}

\begin{exercise}[2]
    用波长为$\lambda$的单色光垂直照射如图所示的劈形膜（$n_3>n_2>n_1$）.观察反射光干涉.葱从劈形膜尖端开始算起，第$2$条明条纹中心所对应的膜厚度为\ans{}.
\end{exercise}

\begin{exercise}[2]{}
    He-Ne激光器发出$\lambda=632.8\mathrm{nm}$的平行光束，垂直照射到一单缝上，在离单缝$2\mathrm{m}$远的屏上观察弗琅禾费衍射图样.测得两个第二级暗纹间的距离是$10\mathrm{cm}$，则单缝的宽度$a=$\ans{}.
\end{exercise}

\begin{exercise}
    用波长为$589.3\mathrm{nm}$钠黄光垂直入射在每毫米有$500$条缝的光栅上，第一级主极大的衍射角是\ans{}.
\end{exercise}

\begin{exercise}[6]{}
    一束光垂直入射在偏振片$P$上，以入射光线为轴转动$P$，观察通过$P$的光强的变化过程.若入射光是\ans{自然}光，则将看到光强不变；若入射光是\ans{线偏振光}光，则将看到明暗交替变化，有时出现全暗；若入射光是\ans{部分偏振光}光，则将看到明暗交替变化.
\end{exercise}

\begin{exercise}[3]{光的偏振}
    使光强为$I_0$的自然光依次垂直通过偏振片$P_1,\ P_2$和$P_3$.$P_1$与$P_2$的偏振化方向成$45^{\circ}$角，$P_2$与$P_3$的偏振化方向成$45^{\circ}$角.则透过三块偏振片的光强$I$为\ans{$I_0/4$}.
\end{exercise}
\solute{由Malus定律得$I=I_0\cos^2{45^{\circ}}\cos^2{45^\circ}=I_0/4$.}

\begin{exercise}[4]{}
    根据量子论，氢原子中核外电子的状态可由四个量子数来确定，其中主量子数$n$可取值为\ans{$\mathbb{N}^+$}，它可决定\ans{电子的能量}.
\end{exercise}

\section{计算题（共48分）\scriptsize 缺28分}
{
\begin{paracol}{2}
\begin{exercise}[10]{}
    图示一平面简谐波在$t=0$时刻的波形图，求
\begin{enumerate}
    \item 该波的波动表达式.
    \item $P$处质点的振动方程.
\end{enumerate}
\end{exercise}
\end{paracol}
\begin{solution}
\begin{enumerate}
    \item 由图可知$t=0$时$O$处质点速度大于零，所以此时其相位为$\varphi=-\frac{\pi}{2}$，其振动方程为
    $$y_0=0.04\cos{\ab(\omega t-\frac{\pi}{2})}$$
    简谐波波长$\lambda=0.4\mathrm{m}$，角频率$\omega=\frac{2\pi u}{\lambda}=0.4\pi$.
\end{enumerate}
\end{solution}

\begin{exercise}[10]{}
    一横波沿绳子传播，其波的表达式为$y=0.05\cos{\ab(100\pi t-2\pi x)}$ (SI).
\begin{enumerate}
    \item 求此波的振幅、波速、频率和波长.
    \item 求绳子上各质点的最大振动速度和最大振动加速.
    \item 求$x_1=0.2\mathrm{m}$处和$x_2=0.7\mathrm{m}$处二质点振动的相位差.
\end{enumerate}
\end{exercise}
\begin{solution}
\begin{enumerate}
    \item 将$y=A\cos{\ab[2\pi\nu\ab(t-\frac{x}{u})+\varphi]}$与所给波函数比对得$A=0.05\mathrm{m},\ \nu=50\mathrm{Hz},\ u=50\mathrm{m/s},\ \lambda=1.0\mathrm{m}$.\point{6}
    \item $v_{\max}=2\pi\nu A=15.7\mathrm{m/s},\ a_{\max}=\ab(2\pi\nu)^2A=4934.8\mathrm{m/s}^2$.\point{4}
\end{enumerate}
\end{solution}
}